\section{Carrier-aware numerical inverse and acceptance gates}
At a prescribed geometry, the finite-source operators supply generalized loads
$F_i(\bm a)=\bm a^\dagger\bm K_i\bm a$ and sampled graph-displacement
amplitudes $b_s(\bm a)=\bm B_s\bm a$, with
$\bm B_s=\bm V_s/(\omega|n_{z,s}|)$. Both maps arise from the same chamber
solution. The matrices $\bm B_s$ have units of seconds when source commands
have units of velocity. With the actual volume constraints already imposed,
let $\bm K_m$ be the capillary--gravity Hessian and $\bm O$ evaluate a surface
perturbation on the fixed apertures. The surrogate
\begin{equation}
 \bm d(\bm a)=\bm O\bm K_m^{-1}
       [\bm F(\bm a)-\bm F_{\rm required}]
\end{equation}
has units of length but is \emph{not} the coupled equilibrium error.
Its inverse stiffness omits the shape derivative of the acoustic load.
A small surrogate residual cannot replace a fresh solution of the full
stationary equations, much less a stability calculation.

For the selected cycle-mean specification, the carrier penalty may be set
to zero while retaining the full quadratic acoustic load and the same source
constraints. Carrier-aware constructions below are optional stricter studies,
not mandatory gates for mean geometry. A small frozen mean residual still
requires the same independent coupled forward and fixed-command convergence
tests; changing the metric does not convert a surrogate into an achieved shape.

\subsection{Coherent source steps with explicit carrier constraints}
A source-step model at $\bm a_k$ uses the exact directional derivative
\begin{equation}
 D F_i(\bm a_k)[\delta\bm a]
       =2\Rea\{\bm a_k^\dagger\bm K_i\delta\bm a\}.
\end{equation}
The carrier trace is already linear. Thus one may minimize an epigraph
variable $t$ subject to convex constraints on the linearized surrogate,
the exact carrier trace and the physical source-amplitude disks:
\begin{align}
 |d_s(\bm a_k)+D d_s(\bm a_k)[\delta\bm a]|&\leq u_j,
       &&s\text{ on face }j,\\
 |\bm B_s(\bm a_k+\delta\bm a)|&\leq v_j,
       &&s\text{ on face }j,\\
 u_j+v_j&\leq t,&
 |a_{k,\ell}+\delta a_\ell|&\leq a_{\ell,\max}.
\end{align}
A trust region bounds the source step; a separately declared annulus penalty
can control the omitted non-optical region. The sum of separate facewise
maxima is conservative relative to the maximum of the pointwise sum, but both
remain sampling diagnostics. Each proposed step is checked against the
\emph{true quadratic} frozen-geometry load before acceptance. The apparatus,
targets and frequency stay fixed during these iterations: this is neither a
physical trajectory nor target continuation. A numerical termination flag is
reported separately from first-order stationarity, feasibility and attainment
of the geometrical criterion.

Source-coordinate conditioning may use the eigenspaces of
$\bm G=\sum_i\bm K_i^\dagger\bm K_i$. Any eigenvalue truncation restricts the
admissible numerical source space and must be declared. It does not establish
the physical rank of the continuous source operator. More emitting regions
can introduce nearly dependent controls rather than improve reachability.

\subsection{Sampled peak majorants}
An alternative smooth-away-from-zero source objective uses unnormalized
finite-vector norms on each aperture. With $\epsilon=10^{-8}\,\mathrm m$,
define $\bm h_j=\bm d_j/\epsilon$, $\bm c_j=\bm b_j/\epsilon$, and
\begin{equation}
 s_j=\|\bm h_j\|_p+\|\bm c_j\|_p,\qquad p\geq2.
\end{equation}
Then $\epsilon s_j\geq\max_s|d_{j,s}|+\max_s|b_{j,s}|$ on the declared
sample sets, even if the two sets differ. No division by the sample count
is made: such a normalization would destroy this majorant. The implemented
least-squares objective uses the two residuals $s_j$ together with separate
annulus and source-effort penalties. This is not exact minimax optimization.
For a nonzero complex vector,
\begin{equation}
 D\|\bm v\|_p[\delta\bm v]
 =\frac{\sum_k |v_k|^{p-2}\Rea(\overline{v_k}\,\delta v_k)}
 {\|\bm v\|_p^{p-1}}.
\end{equation}
At the all-zero vector a zero subgradient is selected. Directional numerical
checks cover both physical and reduced source coordinates. Neither this
majorant nor a large choice of $p$ controls unsampled extrema or repairs
the omitted acoustic shape derivative in $\bm d$.

A complementary, deliberately conservative source restriction chooses an
orthonormal matrix $\bm N$ with
$\|\bm B\bm N\|_2\leq\epsilon_a/(\sqrt{N_s}a_{\max})$.
Commands in its range satisfying $|a_\ell|\leq a_{\max}$ obey
\begin{equation}
 \|\bm B\bm a\|_\infty\leq\|\bm B\bm N\|_2\|\bm a\|_2
 \leq\epsilon_a.
\end{equation}
The operator norm is checked after numerical eigenspace selection.
This is only a sampled frozen-model carrier bound; failure to match the
required traction in this restricted subspace is not a reachability result
for the full source space.

\subsection{A relaxation diagnostic, not a substitute waveform}
In the lifted variable of \cref{eq:lifted-feasibility}, a sampled carrier bound
is linear:
\begin{equation}
 \tr(\bm B_s^\dagger\bm B_s\bm W)\leq\epsilon_a^2.
\end{equation}
The corresponding minimum-carrier semidefinite relaxation can help distinguish
an unsuccessful coherent inverse from a source-space limitation. A higher-rank
covariance is not a coherent command. Inaccurate numerical solutions that
violate positivity, traction equalities or source bounds supply neither a
feasible design nor a verified lower-bound certificate. Bounds derived in a
truncated source space cannot be extrapolated to the full continuous space.

\subsection{Why the maximum spot also requires slope control}
For fixed indices, write the physical Snell map as
$\bm d_t=\eta_n\bm d+(\sqrt Q-\eta_n c)\bm m$, where
$c=\bm d\cdot\bm m$ and $Q=1-\eta_n^2(1-c^2)>0$.
Its first variation is
\begin{align}
 \delta\bm d_t={}&\eta_n\delta\bm d
 +\left(\frac{\eta_n^2c}{\sqrt Q}-\eta_n\right)\delta c\,\bm m
 +(\sqrt Q-\eta_n c)\delta\bm m,\\
 \delta c={}&\delta\bm d\cdot\bm m+\bm d\cdot\delta\bm m.
\end{align}
At a fixed horizontal coordinate of a graph, its upward normal satisfies
\[
 \delta\bm m=
 \frac{(\Id-\bm m\bm m^T)(-\nabla\delta h,0)}
 {\sqrt{1+|\nabla h|^2}}.
\]
Evaluating the normal at a moving intersection adds the change of its
evaluation point. Both effects, and the preceding interface's change to
$\bm d$, must be included in a complete ray derivative. In particular,
small height error without a slope bound does not bound the outgoing angle;
the derivative also becomes ill-conditioned near $Q=0$.

For propagation from $\bm x$ to a fixed detector plane $z=Z$, let
$\tau=(Z-x_z)/d_z>0$ and $\bm y_\perp=\bm x_\perp+\tau\bm d_\perp$.
Then
\[
 \delta\tau=-\frac{\delta x_z+\tau\delta d_z}{d_z},\qquad
 \delta\bm y_\perp=\delta\bm x_\perp+
       \tau\delta\bm d_\perp+\bm d_\perp\delta\tau.
\]
These expressions explain the amplification of short-scale surface errors
into detector displacement. They do not approximate a diffraction image.

A declared regularity proxy augments the frozen height observations with
\[
 \bm d_s(\bm a)=\ell\bm D_r\bm K_m^{-1}
       [\bm F(\bm a)-\bm F_{\rm required}],
 \qquad [\ell]=\mathrm m.
\]
Here $\bm D_r$ evaluates pupil slopes and $\bm d_s$ again has units of length.
For $\ell=10^{-3}\,\mathrm m$, a $10^{-5}$ slope error contributes like a
$10^{-8}\,\mathrm m$ height error. The already-derived source gradient is
composed with this additional linear observation. It does not remove gravity
from $\bm K_m$ or $\bm F_{\rm required}$. Unnormalized $p$-norms may majorize
sampled height and scaled-slope peaks; for a mean-only specification the
carrier contribution is explicitly omitted, while its diagnostic is retained.
Neither the slope scale nor this majorant is an exact optical bound.

The independent joint gate requires each clear-pupil mean height error to be
at most $10^{-8}\,\mathrm m$ and the maximum geometric radius about the
prescribed optical axis at the fixed detector to be at most
$10^{-6}\,\mathrm m$. RMS radius, centroid recentering, refocusing, and
discarding failed rays cannot substitute for this maximum. The intended
illumination and pupils must be declared and checked for full transmission.
Discrete ray samples must be refined and distinguished from a continuous-pupil
certificate; these tests still do not establish maintained liquid accuracy.

\subsection{Static acoustic feedback in the precision inverse}
The source conditions above must not confuse the mechanical compliance with
the coupled static response. At a fixed reference shape and command, write
\[
 \bm R(\bm q,\bm a)=\bm F_m(\bm q)-\bm F_a(\bm q,\bm a),\quad
 \bm A_q=D_q\bm F_a,\quad
 \bm J=\bm K_m-\bm A_q.
\]
The first-order equilibrium correction solves
$\bm J\delta\bm q=-\bm R$. If $\bm J$ is invertible, the local equilibrium
source derivative is $D_a\bm q=\bm J^{-1}D_a\bm F_a$.
An observation of the actual coupled response therefore contains
$\bm O\bm J^{-1}$, and the fixed-detector optical observation contains
$\bm J_{\rm spot}\bm J^{-1}$. The latter includes both surfaces, not
separate single-face focus errors. A small
$\bm O\bm K_m^{-1}\bm R$ need not imply a small
$\bm O\bm J^{-1}\bm R$: the acoustic shape feedback can strongly amplify
some force defects, including those in the non-optical annuli.

As an explicitly truncated diagnostic, choose admissible shape directions
in the columns of $\bm V$ and, for positive definite $\bm K_m$, define
\[
 \bm W=(\bm V^T\bm K_m\bm V)^{-1}\bm V^T\bm K_m,
 \qquad \bm L=\bm K_m^{-1}\bm A_q\bm V.
\]
Replacing $\bm A_q$ by $(\bm A_q\bm V)\bm W$ gives
\begin{equation}
 \bm C_V=\bm K_m^{-1}
 +\bm L(\Id-\bm W\bm L)^{-1}\bm W\bm K_m^{-1}.
\end{equation}
Direct multiplication verifies this inverse whenever the small matrix is
nonsingular. Central differences of newly solved wave fields can measure
$\bm A_q\bm V$; the step and retained directions must be refined.
This is a restriction of the general coupled linearization, not a replacement
for it or for the continuous-source inverse conditions.

Holding $\bm C_V$ fixed during a source step supplies only a local model.
The acoustic derivative itself depends on the command, so finite source
changes require fresh coupled validation and, when needed, an updated
response operator. Rank-truncated static conditioning is neither a stability
spectrum nor an infeasibility certificate. In particular it does not test
three-dimensional disturbances, inertia, dissipative dynamics, a streaming
base flow or thermal feedback.

\subsection{Required independent verification}
For any candidate, solve the nonlinear coupled equilibrium with fixed physical
commands and refine both wave and surface spaces. The absolute mean error and
carrier amplitude must be evaluated independently of design collocation.
Passing a sampled ten-nanometre check still requires a numerical-error budget
and control of unsampled extrema. Formation, stability about the appropriate
streaming base flow, thermal feedback, material uncertainty and experimental
validation remain additional gates. The numerical inverse does not remove
those requirements by using more emitters or higher carrier frequencies.
